\paragraph{Reproduction and limitations.}
In the local reproduction, the Lee and Lee pipeline was implemented as a leakage-aware adaptation of the original feature-selection and classification workflow. The reproduced components include train-only DEG selection with limma, a VAE or autoencoder-style latent representation fitted only on the training data, transcription-factor and CFG-derived subsets obtained by intersecting the train-selected DEGs with the official supplementary gene lists, and classical classifiers trained on the resulting feature spaces. The hub-gene representation was not reproduced in the strict benchmark, because the complete HPRD gene-gene interaction network required to recompute hub genes was not available in the official supplementary material. Therefore, the strongest comparable local configuration is the DEG-based representation combined with regularized classical classifiers.

The main limitation of the original Lee and Lee study is that the predictive task is AD versus cognitively normal controls, while MCI is not the main classification target. The original paper is valuable because it includes external validation across datasets, but its results still do not directly answer the AD versus MCI question addressed in this thesis. A second limitation is that the original workflow contains several dataset-level preprocessing and harmonization steps, including cross-dataset batch correction. Such steps are reasonable for building a common expression space, but they must be handled carefully in a benchmark setting, because estimating normalization, batch correction, or feature-selection parameters before the validation split can make the evaluation optimistic. Finally, part of the original biological feature engineering depends on external resources and supplementary files that are not always sufficient to reconstruct every feature set exactly. For this reason, the reproduction emphasizes the components that can be implemented transparently and fitted inside each training fold.

\paragraph{Reproduction and limitations.}
In the local reproduction, the Kelly et al. approach was adapted to the local AD versus MCI and AD versus CTL benchmark by reproducing the main feature-selection and classifier families of the paper. The implementation includes LASSO selection, VSSRFE with logistic regression, VAE-based dimensionality reduction, and the classical classifiers used in the original study, with particular attention to the VSSRFE-LR plus random forest configuration. The paper-like 159-gene VSSRFE setting and the corresponding logistic-regression regularization value were also considered when a fixed-paper configuration was required. All preprocessing, scaling, feature selection, VAE fitting, and model tuning were moved inside the training split in the strict benchmark.

The main limitation of the original Kelly et al. AD analysis is task mismatch. MCI samples are explicitly removed to reduce complexity, so the reported AD model is a disease-versus-control classifier rather than an AD-versus-MCI classifier. This makes the reported ROC-AUC a useful reference for blood-expression classification, but not a direct estimate of early-stage discrimination. A second limitation is that the original AD setting uses GSE63061 and GSE63060 as independent training and test datasets from the AddNeuroMed study. This is stronger than simple pooled cross-validation, but it still evaluates a setting with related cohorts and without the MCI boundary. In addition, the best model depends on an expensive wrapper feature-selection procedure. In high-dimensional transcriptomic data, such procedures can become sensitive to the exact training split and to the metric used during internal optimization. The reproduction therefore treats VSSRFE as a strong but potentially unstable classical baseline rather than as a definitive biomarker panel.

\paragraph{Reproduction and limitations.}
In the local reproduction, One2MFusion was implemented by retaining the central structure of the original method: LASSO-based gene selection, Fisher-score-based gene grouping, supervised construction of a two-dimensional gene-expression image through an LDA-style mapping, a feed-forward branch for the selected gene vector, a convolutional branch for the generated image, and late fusion of the two neural representations. The reproduction also includes the paper-like fixed-gene-count setting for AD versus MCI, while enforcing that the imputer, scaler, LASSO selection, Fisher grouping, image mapping, and neural-model selection are fitted only on the training portion of each split.

The original One2MFusion paper is important because it directly includes AD, MCI, and control comparisons, but it is also the method with the clearest risk of optimistic evaluation. The gene selection, Fisher-distance grouping, and LDA image construction are supervised transformations. If they are fitted on the full dataset before cross-validation, information from the validation fold can influence the feature space and the image representation even before the classifier is trained. This is particularly problematic because the method reports five-fold cross-validation on a pooled dataset rather than an independent external validation. A second limitation is that the image representation imposes an artificial spatial structure on gene expression data. Convolutional filters may exploit this constructed layout, but the biological meaning of local image neighborhoods is not guaranteed. The local reproduction therefore evaluates One2MFusion mainly as a representation-learning hypothesis, while treating the original reported performance as potentially optimistic unless all supervised transformations are nested inside the training folds.

\paragraph{Reproduction and limitations.}
In the local reproduction, the Diagnostics 2025 gene-expression branch was adapted to the transcriptome-only binary benchmark. The reproduced pipeline includes train-only MinMax scaling, XGBoost-based gene ranking, selection of a reduced subset through an approximate SFBS procedure, optional Borderline-SMOTE applied only to the training data, and downstream classifiers including the dense deep-learning model used as the main comparison point. The clinical-data branch was not reproduced, because the thesis benchmark is restricted to blood gene expression and because importing clinical variables would make the comparison with transcriptome-only methods unfair.

The original Diagnostics 2025 study has two important limitations for the present thesis. First, it is a multiclass ADNI stage-diagnosis study, not a binary AD versus MCI study on GSE63060 and GSE63061. The reported CN, MCI, and dementia scores therefore cannot be compared directly with the local pairwise AD/MCI results. Second, the paper shows that clinical variables perform substantially better than gene expression, which is clinically informative but also confirms that transcriptome-only prediction remains limited. A further limitation is that the published gene-expression branch selects a compact set of probes from an HDLSS and imbalanced dataset. This makes the result sensitive to feature-ranking stability and to the precise handling of class imbalance. In the reproduction, oversampling is therefore restricted to the training split, and original validation and test samples are preserved to avoid inflating performance with synthetic observations.

\paragraph{Reproduction and limitations.}
In the local reproduction, the Scientific Reports 2026 pipeline was adapted by implementing the feature-selection and neural-classification components that are compatible with the common benchmark. The reproduced selectors include Chi-square, ANOVA, RFE, and ElasticNet-style ranking, with ElasticNet plus a dense neural network used as the most relevant reproduced configuration. The local implementation also includes train-only tabular GAN augmentation as an approximation of the paper's TGAN/CTGAN setup. Synthetic samples are generated only from training data after preprocessing and feature selection. SHAP ranking and GO/KEGG enrichment were not included in the reproduced benchmark, because they are interpretability analyses rather than core predictive components and would not change the common evaluation protocol.

The main limitation of the original Scientific Reports 2026 paper is again task mismatch: it focuses on AD versus healthy control classification using GSE63060, GSE63061, ADNI, and integrated datasets, rather than on AD versus MCI. The strong reported accuracy therefore reflects a more separable diagnostic contrast. A second limitation is dataset integration. The paper acknowledges that the datasets come from different experimental platforms and that formal batch-effect correction such as ComBat was not applied. This is important because integrated transcriptomic datasets can contain technical structure that is easier for a model to exploit than disease biology. Finally, GAN-based augmentation is attractive in small-sample transcriptomics, but it can also amplify artifacts from the training distribution or create overly smooth synthetic samples. For this reason, the reproduction treats augmentation as an ablation inside a leakage-aware pipeline, not as an automatic solution to the limited-sample problem.
